% ============================================================================
% general-character.tex  (v3, 2026-08-16)
%
% Restructures Section 2 ("An interpolation theorem") of
% paper/weighted-races.tex (the 2829-line draft of 2026-08-15/16) so that
% Theorem~\ref{thm:interp} is stated and proved for a GENERAL real primitive
% Dirichlet character chi mod q, with chi_4 and chi_3 as worked examples ---
% the referee's point: Theorem~\ref{thm:dissolution} is already general,
% and the generality of Theorem~\ref{thm:interp} is free.  Also fixes, in
% the same edit set, the stale sentence in the Scope paragraph of the proof
% of (iii) ("...Theorem~\ref{thm:dissolution} of the companion draft and is
% \emph{not} proved in this paper" -- the dissolution law IS proved in this
% paper, in Section~\ref{sec:dissolution}), and the follow-on stale scope
% statements in Sections 4-5.
%
% FORMAT.  Numbered string-anchored replacement blocks against the CURRENT
% weighted-races.tex.  Each block gives the EXACT current text (OLD, every
% line prefixed "%OLD% "; blank lines as bare "%OLD%"; original line breaks
% preserved) and the exact replacement (NEW, between "%---NEW--->" and
% "%<---NEW---", uncommented).  Every OLD snippet was machine-verified to
% occur exactly once in the current file; the full edit set was applied to
% a scratch copy and compiled with pdflatex (all references resolved).
%
% MATHEMATICAL STRUCTURE OF THE GENERALIZATION.
%  * Hypotheses: GRH(chi) + NCZ(chi): L(1/2,chi) != 0, with LI(chi) where
%    densities are asserted.  NCZ is unconditional for q = 3, 4 by grouped
%    alternating series (chi_4: alternating series test; chi_3: blocks of
%    three, >= 1 - 1/sqrt(2) > 0).  gamma != 0 is DERIVED from GRH + NCZ
%    (a zero with gamma = 0 would sit at s = 1/2), never built into LI --
%    same structure as the current chi_4 text.  gamma_1(chi) > 0 replaces
%    every use of the numeric gamma_1 > 6 (which held for chi_4); constants
%    become chi-dependent, legitimate since the theorem is for fixed chi.
%  * Explicit formula (the crux): for odd chi (a=1) the proof is as before
%    (L(0,chi) != 0 by the functional equation and L(1,chi) != 0).  For
%    even chi (a=0), s=0 is a simple TRIVIAL zero: L'/L(s) = 1/s + b(chi)
%    + O(s), the Perron kernel's 1/s pole meets it, and the s=0 residue of
%    (-L'/L)(s) x^s/s is -log x - b(chi) (Davenport Ch. 19's -(1-a) log x
%    - b(chi)); the left contour edge moves to U = 2K+1-a so it passes
%    between trivial zeros (at -(2j+a)) in either parity.  This changes
%    ONLY the O(log x) floor of eq:classical, which the downstream proofs
%    already absorb (verified: the floor enters lem:trunc as (1+x^m) log x
%    and the ANS mean-square hypothesis through u^2 e^{-kappa_m u}).
%  * Prime squares: sum_{p^2<=x} chi(p)^2 p^{2m} log p =
%    sum_{p<=sqrt x, p coprime q} p^{2m} log p; the at most omega(q) primes
%    p | q each contribute p^{2m} log p <= (1+q^{2m}) log q = O_{q,m}(1),
%    absorbed (Block 15).  The Methods paragraph's earlier bound
%    "O_m(x^m log x)" for these terms was WRONG for -1/2 < m < 0 (there
%    x^m log x -> 0 while the terms are fixed positive constants); Block 34
%    corrects it to O_{q,m}(1).
%  * C(chi) := sum_{gamma>0} 2/gamma^2 = (log Lambda)''(1/2, chi), with
%    Lambda(s,chi) = (q/pi)^{(s+a)/2} Gamma((s+a)/2) L(s,chi) and root
%    number +1 for every real primitive chi (Gauss).
%
% NUMERICS (all recomputed for this edit set with mpmath at 50 digits;
% zero data from prototype/explore/chi{3,4}_zeros.txt, 537/511 ordinates):
%  * C(chi_4) = (1/4)psi'(3/4) + (log L)''(1/2,chi_4)
%             = 0.15602964988964488...  (matches the paper; psi'(3/4)
%             = pi^2 - 8G checked).  Zero-sum cross-check 0.152546 +
%             0.003487 = 0.156033 reproduced exactly.
%  * C(chi_3) = (1/4)psi'(3/4) + (log L)''(1/2,chi_3)
%             = 0.11340215685275857... (matches the certified value
%             recorded in Section 5); (log L)''(1/2,chi_3)
%             = -0.52206775506514...  NEW zero-sum cross-check computed
%             for Block 3: 537 ordinates give 0.110308, RvM tail
%             (dN = (2 pi)^{-1} log(3t/(2 pi)) dt from gamma_537 =
%             699.7747...) gives 0.003098, total 0.113406, i.e. closed
%             form + 4.1e-6 -- same O(1e-5) fluctuation scale as chi_4.
%  * Thresholds (truncated digits): chi_4: m_0 = 0.5750748... ("0.57507"),
%    Chebyshev < 1/2 for m < 0.3950585... (the OLD text printed the
%    ROUNDED "0.39506\ldots"; Blocks 4/23 print the truncation
%    "0.395058\ldots"), Chebyshev >= 1 from 0.7658039... ("0.76580").
%    chi_3: m_0 = 0.7610480... ("0.761048"), Chebyshev < 1/2 for
%    m < 0.5498913... ("0.549891"), >= 1 from 0.9847705... ("0.984770");
%    exp(-1/(2 C(chi_3))) = 0.0121662... ("0.01216"); 1 - delta_3(0) =
%    1 - 0.999063 = 9.37e-4.
%  * L(1/2,chi_4) = 0.667691457..., L(1/2,chi_3) = 0.4808675...,
%    gamma_1(chi_4) = 6.0209..., gamma_1(chi_3) = 8.0397... (all match
%    the paper); delta_3(0) = 0.999063 and the mod-3 first-flip facts
%    (p=13 at m=1, p=7 at m=2,3; no classical flip below 1e10) quoted
%    verbatim from the introduction.
%
% WHAT IS DELIBERATELY NOT CHANGED (honesty ledger):
%  * Part (iv) of thm:interp stays a chi_4 statement: the endpoint results
%    (AK23/Sheth24/Hay25) are quoted by this paper for the mod-4 race.
%  * Corollary~\ref{cor:dissolution-race} stays q=4; the general race
%    transfer is stated in the scope remarks (Blocks 23, 32, 35) as an
%    immediate combination of thm:dissolution + thm:interp(ii) under
%    GRH + LI + NCZ for the given chi.
%  * Section 3 (variance in closed form) remains the q=4 computation, as
%    flagged in Block 30; its general-q content is already
%    Proposition~\ref{prop:varasymp} / eq:varidq.
%  * For general q, NCZ remains an assumption (Chowla's conjecture); the
%    paper proves it only for q = 3, 4.
%
% NEW LABEL: cor:mod3rate (Block 4).  No labels are removed.  The new
% Corollary changes the printed numbering of later corollaries (the
% corollary counter is independent in this preamble); all \ref's resolve.
%
% REFEREE ADDENDA (Blocks 36-39, added on verification).  Four stale
% lines the original 35 blocks left behind, found by refereeing the
% fully patched file:
%  * Block 36: the lem:trunc proof opened "Since psi(t,chi)=0 for
%    t<3" -- TRUE for chi_4 (chi_4(2)=0) but FALSE for any chi with
%    2 coprime to q (e.g. chi_3: psi(2,chi_3) = -log 2 != 0).  The
%    displayed partial-summation identity is still exact (the lower
%    limit 2^- captures the jump at 2 and psi(2^-,chi)=0), so only the
%    justification sentence needed repair: t<3 -> t<2.
%  * Block 37: Step 2 of the lem:classical proof said "#G <= A_1
%    log(T+2) for an absolute A_1"; A_1 comes from (F3), which is now
%    <<_q, and Block 7 declares all constants q-dependent -- so
%    "absolute" was an internal contradiction: now A_1 = A_1(q).
%  * Blocks 38-39: the abstract's item (iii) and Results item 3 still
%    said the LI transfer of the dissolution law applies "to the
%    mod-4 race itself" -- true but stale understatements once
%    thm:interp(ii) is general (the body's Blocks 23/31/32/35 already
%    state the general transfer); both now say it holds for every
%    real primitive chi under LI(chi), with cor:dissolution-race
%    recording q=4.
% ============================================================================

% ============================================================================
% Abstract: the interpolation-theorem sentence now reflects the general
% statement (arbitrary real primitive chi; C(chi) closed form; chi_4 and
% chi_3 values).  Lines 60-67.
% ============================================================================
%% BLOCK 1 REPLACES:
%OLD% theorem, proved by the framework of Akbary--Ng--Shahabi and the method
%OLD% of Humphries: under GRH and LI for $L(s,\chi_4)$, for every
%OLD% $m\in(-\tfrac12,\infty)$ the logarithmic density $\delta(m)$ of the set
%OLD% of $x$ where the class $3\ (\mathrm{mod}\ 4)$ leads exists and lies in
%OLD% $(\tfrac12,1)$, the map $m\mapsto\delta(m)$ is continuous, and
%OLD% $1-\delta(m)\le\exp\bigl(-1/(2C(2m+1)^2)\bigr)$ with
%OLD% $C=\sum_{\gamma>0}2\gamma^{-2}=0.1560296\ldots$ (evaluated in closed
%OLD% form), so $\delta(m)\to1$ as $m\to-\tfrac12^{\,+}$. This exhibits the classical density formulation
%---NEW--->
theorem, proved by the framework of Akbary--Ng--Shahabi and the method
of Humphries and stated for an arbitrary real primitive character
$\chi$ mod $q$: under GRH$(\chi)$, LI$(\chi)$ and
$L(\tfrac12,\chi)\ne0$ (the nonvanishing is unconditional for
$q=3,4$), for every
$m\in(-\tfrac12,\infty)$ the logarithmic density $\delta_\chi(m)$ of
the set of $x$ where the weighted nonresidue count leads (for
$\chi=\chi_4$, where the class $3\ (\mathrm{mod}\ 4)$ leads) exists
and lies in
$(\tfrac12,1)$, the map $m\mapsto\delta_\chi(m)$ is continuous, and
$1-\delta_\chi(m)\le\exp\bigl(-1/(2C(\chi)(2m+1)^2)\bigr)$ with
$C(\chi)=\sum_{\gamma>0}2\gamma^{-2}$ evaluated in closed form
($C(\chi_4)=0.1560296\ldots$, $C(\chi_3)=0.1134021\ldots$), so
$\delta_\chi(m)\to1$ as $m\to-\tfrac12^{\,+}$. This exhibits the classical density formulation
%<---NEW---

% ============================================================================
% Results item 5: state the theorem's actual (general) hypotheses and
% constants.  Lines 284-298.
% ============================================================================
%% BLOCK 2 REPLACES:
%OLD% \textbf{5. An interpolation theorem.} Section~\ref{sec:theorem}
%OLD% (Theorem~\ref{thm:interp}) proves: under GRH and LI for $L(s,\chi_4)$,
%OLD% for every $m\in(-\tfrac12,\infty)$ the
%OLD% normalized race has a limiting logarithmic distribution, the density
%OLD% $\delta(m)$ exists and lies in $(\tfrac12,1)$, the map
%OLD% $m\mapsto\delta(m)$ is continuous, and
%OLD% \[
%OLD% 1-\delta(m)\;\le\;\min\Bigl\{C\,(2m+1)^2,\;
%OLD% e^{-1/(2C(2m+1)^2)}\Bigr\},\qquad
%OLD% C=\sum_{\gamma>0}\frac{2}{\gamma^2}=0.1560296\ldots,
%OLD% \]
%OLD% so $\delta(m)\to1$ as $m\to-\tfrac12^{\,+}$ --- a weight-aspect
%OLD% instance of the highly biased phenomenon of Fiorilli~\cite{Fio14} ---
%OLD% matching the single-sign
%OLD% regime established at $m=-\tfrac12$ by \cite{AK23,Sheth24,Hay25}.
%---NEW--->
\textbf{5. An interpolation theorem.} Section~\ref{sec:theorem}
(Theorem~\ref{thm:interp}) proves, for an arbitrary real primitive
character $\chi$ mod $q$ under GRH$(\chi)$, LI$(\chi)$ and
$L(\tfrac12,\chi)\ne0$ (the nonvanishing is unconditional for
$q=3,4$): for every $m\in(-\tfrac12,\infty)$ the
normalized race has a limiting logarithmic distribution, the density
$\delta_\chi(m)$ exists and lies in $(\tfrac12,1)$, the map
$m\mapsto\delta_\chi(m)$ is continuous, and
\[
1-\delta_\chi(m)\;\le\;\min\Bigl\{C(\chi)\,(2m+1)^2,\;
e^{-1/(2C(\chi)(2m+1)^2)}\Bigr\},\qquad
C(\chi)=\sum_{\gamma>0}\frac{2}{\gamma^2},
\]
with $C(\chi_4)=0.1560296\ldots$ and $C(\chi_3)=0.1134021\ldots$ (both
evaluated in closed form),
so $\delta_\chi(m)\to1$ as $m\to-\tfrac12^{\,+}$ --- a weight-aspect
instance of the highly biased phenomenon of Fiorilli~\cite{Fio14} ---
matching, at $q=4$, the single-sign
regime established at $m=-\tfrac12$ by \cite{AK23,Sheth24,Hay25}.
%<---NEW---

% ============================================================================
% Section 2 preamble: general real primitive chi, NCZ hypothesis and
% gamma_1(chi)>0, general R_{m,chi}/E_{m,chi}/X_m, C(chi) closed form;
% chi_4 and chi_3 as worked examples.  Lines 312-365.
% ============================================================================
%% BLOCK 3 REPLACES:
%OLD% Throughout this section $\chi=\chi_4$ is the nontrivial character mod
%OLD% $4$, $\rho=\tfrac12+i\gamma$ runs over the nontrivial zeros of
%OLD% $L(s,\chi)$ (listed with multiplicity), GRH$(\chi)$ denotes the Riemann
%OLD% hypothesis for $L(s,\chi)$, and LI$(\chi)$ the linear independence over
%OLD% $\mathbb Q$ of the multiset of positive ordinates $\gamma$ (LI in
%OLD% particular forces all $\gamma$ to be distinct, i.e.\ the zeros to be
%OLD% simple; recall also that $L(\tfrac12,\chi_4)=\sum_{k\ge0}(-1)^k(2k+1)^{-1/2}>0$
%OLD% by the alternating series test, so no ordinate vanishes). For
%OLD% $m\in(-\tfrac12,\infty)$ define the ($\theta$-form) weighted race
%OLD% \[
%OLD% R_m(x)\;=\;\sum_{\substack{p\le x\\ p\equiv3\,(4)}}p^m\log p
%OLD% \;-\;\sum_{\substack{p\le x\\ p\equiv1\,(4)}}p^m\log p
%OLD% \;=\;-\sum_{p\le x}\chi(p)\,p^m\log p ,
%OLD% \]
%OLD% and its normalization
%OLD% \[
%OLD% E_m(x)\;=\;(2m+1)\,x^{-(m+1/2)}\,R_m(x).
%OLD% \]
%OLD% Write
%OLD% \[
%OLD% a_\gamma(m)\;=\;\frac{2m+1}{\sqrt{(m+\tfrac12)^2+\gamma^2}},
%OLD% \qquad
%OLD% X_m\;=\;1+\sum_{\gamma>0}2a_\gamma(m)\cos\theta_\gamma ,
%OLD% \]
%OLD% with $\theta_\gamma$ i.i.d.\ uniform on $[0,2\pi)$; the series converges
%OLD% a.s.\ and in $L^2$ since $\sum_\gamma a_\gamma(m)^2<\infty$ (the ordinates
%OLD% have density $\ll\log\gamma$). Set
%OLD% \[
%OLD% C\;:=\;\sum_{\gamma>0}\frac{2}{\gamma^2}\;=\;0.1560296498896\ldots
%OLD% \]
%OLD% Under GRH$(\chi_4)$ (in force throughout this section) this constant has
%OLD% a closed evaluation: the completed function
%OLD% $\Lambda(s)=(4/\pi)^{(s+1)/2}\Gamma\bigl(\tfrac{s+1}{2}\bigr)L(s,\chi_4)$
%OLD% is entire and satisfies $\Lambda(s)=\Lambda(1-s)$ (the root number of
%OLD% $\chi_4$ is $+1$), so $\Xi(z):=\Lambda(\tfrac12+z)$ is even; since
%OLD% $\Xi(0)>0$ (as $L(\tfrac12,\chi_4)>0$), $\sum_\gamma\gamma^{-2}<\infty$,
%OLD% and, under GRH, the zeros of $\Xi$ are exactly $\pm i\gamma$, the
%OLD% Hadamard factorization takes the genus-zero form
%OLD% $\Xi(z)/\Xi(0)=\prod_{\gamma>0}(1+z^2/\gamma^2)$, exactly as for the
%OLD% Riemann $\xi$-function \cite[\S12]{Dav00}. Hence
%OLD% \[
%OLD% C\;=\;(\log\Lambda)''\bigl(\tfrac12\bigr)
%OLD% \;=\;\tfrac14\,\psi'\bigl(\tfrac34\bigr)
%OLD% +\frac{d^2}{ds^2}\log L(s,\chi_4)\Big|_{s=1/2}
%OLD% \;=\;0.15602964988964488\ldots,
%OLD% \]
%OLD% with $\psi'(\tfrac34)=\pi^2-8G$ ($G$ Catalan's constant), evaluated at
%OLD% $40$-digit working precision. Cross-check from the zero data: the first
%OLD% $511$ ordinates give $0.152546$ and the Riemann--von Mangoldt tail
%OLD% $\int_{\gamma_{511}}^\infty 2t^{-2}\,dN(t)$ with
%OLD% $dN(t)=(2\pi)^{-1}\log(2t/\pi)\,dt$ gives $0.003487$, total $0.156033$,
%OLD% agreeing with the exact value to $3.4\times10^{-6}$ --- within the
%OLD% expected $O(10^{-5})$ fluctuation of the zero-counting remainder.
%OLD% (Earlier drafts printed the cruder value $0.156033$.)
%---NEW--->
Throughout this section $\chi$ is a \emph{fixed} real primitive
Dirichlet character mod $q\ge3$, of parity $\chi(-1)=(-1)^a$ with
$a\in\{0,1\}$; the two moduli computed in this paper are $\chi_4$ (the
nontrivial character mod $4$; $a=1$) and $\chi_3$ (the quadratic
character mod $3$; $a=1$), and they serve as running examples.
$\rho=\beta+i\gamma$ runs over the nontrivial zeros of
$L(s,\chi)$ (listed with multiplicity), GRH$(\chi)$ denotes the Riemann
hypothesis for $L(s,\chi)$, and LI$(\chi)$ the linear independence over
$\mathbb Q$ of the multiset of positive ordinates $\gamma$ (LI in
particular forces all $\gamma$ to be distinct, i.e.\ the zeros to be
simple; we do \emph{not} build the nonvanishing of ordinates into the
LI convention). Alongside GRH$(\chi)$ we assume throughout the
no-central-zero hypothesis
\[
\mathrm{NCZ}(\chi):\qquad L(\tfrac12,\chi)\neq0 ,
\]
which under GRH forces every ordinate to be nonzero (a zero with
$\gamma=0$ would sit at $s=\tfrac12$), so that the least positive
ordinate $\gamma_1(\chi)$ is a strictly positive number --- possibly
small for large $q$; the theorem is for fixed $\chi$, and all implied
constants in this section may depend on $\chi$ (equivalently on $q$)
without further mention. For $q=3,4$ the hypothesis NCZ is
unconditional, by grouped alternating series: for $\chi_4$,
$L(\tfrac12,\chi_4)=\sum_{k\ge0}(-1)^k(2k+1)^{-1/2}>0$ by the
alternating series test, while for $\chi_3$ the Dirichlet series is
not alternating (the vanishing terms $\chi_3(3k)=0$ interleave), but
grouping it in blocks of three gives
$L(\tfrac12,\chi_3)=\sum_{k\ge0}\bigl[(3k+1)^{-1/2}-(3k+2)^{-1/2}\bigr]
\ge1-2^{-1/2}>0$ (justification and numerics in
Section~\ref{sec:methods}: $L(\tfrac12,\chi_4)=0.667691457\ldots$,
$L(\tfrac12,\chi_3)=0.4808675\ldots$). For general real primitive
$\chi$, NCZ is central nonvanishing --- Chowla's conjecture, known in
every numerically checked case but open in general; under GRH,
$L(\tfrac12,\chi)\ge0$, so NCZ is then exactly $L(\tfrac12,\chi)>0$
(see the discussion in Section~\ref{sec:methods}). For
$m\in(-\tfrac12,\infty)$ define the ($\theta$-form) weighted race
\[
R_{m,\chi}(x)\;=\;-\sum_{p\le x}\chi(p)\,p^m\log p
\;=\;\sum_{\substack{p\le x\\ \chi(p)=-1}}p^m\log p
\;-\;\sum_{\substack{p\le x\\ \chi(p)=+1}}p^m\log p
\]
(nonresidues minus residues; for $\chi=\chi_4$ this is the race of
$p\equiv3$ against $p\equiv1$ mod $4$, for $\chi_3$ of $p\equiv2$
against $p\equiv1$ mod $3$), and its normalization
\[
E_{m,\chi}(x)\;=\;(2m+1)\,x^{-(m+1/2)}\,R_{m,\chi}(x).
\]
Since $\chi$ is fixed we abbreviate $R_m=R_{m,\chi}$ and
$E_m=E_{m,\chi}$. Write
\[
a_\gamma(m)\;=\;\frac{2m+1}{\sqrt{(m+\tfrac12)^2+\gamma^2}},
\qquad
X_m\;=\;1+\sum_{\gamma>0}2a_\gamma(m)\cos\theta_\gamma ,
\]
with $\theta_\gamma$ i.i.d.\ uniform on $[0,2\pi)$; the series converges
a.s.\ and in $L^2$ since $\sum_\gamma a_\gamma(m)^2<\infty$ (the
ordinates have density $\ll_q\log\gamma$). ($X_m$ is the $X_{m,\chi}$
of Section~\ref{sec:dissolution}.) Set
\[
C(\chi)\;:=\;\sum_{\gamma>0}\frac{2}{\gamma^2}\;<\;\infty .
\]
Under GRH$(\chi)$ and NCZ$(\chi)$ (in force throughout this section)
this constant has a closed evaluation: the completed function
$\Lambda(s,\chi)=(q/\pi)^{(s+a)/2}\Gamma\bigl(\tfrac{s+a}{2}\bigr)
L(s,\chi)$
is entire and satisfies $\Lambda(s,\chi)=\Lambda(1-s,\chi)$ --- every
real primitive character has root number $+1$, by Gauss's evaluation
of the sign of the quadratic Gauss sum, $\tau(\chi)=\sqrt q$ for
$a=0$ and $i\sqrt q$ for $a=1$ (\cite[Chs.~2, 9]{Dav00}; see the proof
of Proposition~\ref{prop:varasymp}) --- so
$\Xi(z):=\Lambda(\tfrac12+z,\chi)$ is even; since $\Xi(0)\ne0$ (by
NCZ; under GRH in fact $\Xi(0)>0$),
$\sum_\gamma\gamma^{-2}<\infty$,
and, under GRH, the zeros of $\Xi$ are exactly $\pm i\gamma$, the
Hadamard factorization takes the genus-zero form
$\Xi(z)/\Xi(0)=\prod_{\gamma>0}(1+z^2/\gamma^2)$, exactly as for the
Riemann $\xi$-function \cite[\S12]{Dav00} (the argument is written
out, at $q=4$, in Proposition~\ref{prop:varid}, and applies verbatim).
Hence
\[
C(\chi)\;=\;(\log\Lambda)''\bigl(\tfrac12,\chi\bigr)
\;=\;\tfrac14\,\psi'\Bigl(\frac{2a+1}{4}\Bigr)
+\frac{d^2}{ds^2}\log L(s,\chi)\Big|_{s=1/2}.
\]

\medskip\noindent\emph{Example ($q=4$).} For $\chi=\chi_4$ ($a=1$;
$\gamma_1(\chi_4)=6.0209\ldots$),
\[
C(\chi_4)\;=\;\tfrac14\,\psi'\bigl(\tfrac34\bigr)
+\frac{d^2}{ds^2}\log L(s,\chi_4)\Big|_{s=1/2}
\;=\;0.15602964988964488\ldots,
\]
with $\psi'(\tfrac34)=\pi^2-8G$ ($G$ Catalan's constant), evaluated at
$40$-digit working precision (Corollary~\ref{cor:Cexact}). Cross-check
from the zero data: the first
$511$ ordinates give $0.152546$ and the Riemann--von Mangoldt tail
$\int_{\gamma_{511}}^\infty 2t^{-2}\,dN(t)$ with
$dN(t)=(2\pi)^{-1}\log(2t/\pi)\,dt$ gives $0.003487$, total $0.156033$,
agreeing with the exact value to $3.4\times10^{-6}$ --- within the
expected $O(10^{-5})$ fluctuation of the zero-counting remainder.
(Earlier drafts printed the cruder value $0.156033$.)

\medskip\noindent\emph{Example ($q=3$).} For $\chi=\chi_3$ ($a=1$;
$\gamma_1(\chi_3)=8.0397\ldots$, from the $537$ ordinates computed to
$25$ digits in Section~\ref{sec:methods}),
\[
C(\chi_3)\;=\;\tfrac14\,\psi'\bigl(\tfrac34\bigr)
+\frac{d^2}{ds^2}\log L(s,\chi_3)\Big|_{s=1/2}
\;=\;0.11340215685275857\ldots,
\]
the second derivative being
$(\log L)''(\tfrac12,\chi_3)=-0.522067755\ldots$; the value is
computed by the same closed-form route as $C(\chi_4)$ and is recorded
among the certified constants of Section~\ref{sec:methods}.
Cross-check from the zero data: the $537$ ordinates give $0.110308$
and the tail $\int_{\gamma_{537}}^\infty 2t^{-2}\,dN(t)$ with
$dN(t)=(2\pi)^{-1}\log\bigl(3t/(2\pi)\bigr)\,dt$ gives $0.003098$,
total $0.113406$, agreeing with the closed form to
$4.1\times10^{-6}$.
%<---NEW---

% ============================================================================
% Theorem thm:interp restated for a general real primitive chi
% (hypotheses GRH(chi) + NCZ(chi), parts (i)-(iii) under LI(chi));
% part (iv) is the chi_4 anchor to the endpoint literature, whose
% citations (AK23, Sheth24, Hay25) this paper quotes for the mod-4 race
% only.  Numeric thresholds move to the examples; new mod-3 rate
% corollary (task (d)).  Lines 367-415.
% ============================================================================
%% BLOCK 4 REPLACES:
%OLD% \begin{theorem}\label{thm:interp}
%OLD% Assume GRH$(\chi_4)$.
%OLD% \begin{enumerate}
%OLD% \item[(i)] For each $m\in(-\tfrac12,\infty)$ the function
%OLD% $u\mapsto E_m(e^u)$ possesses a limiting distribution
%OLD% $\nu_m$: for every bounded continuous $f:\mathbb R\to\mathbb R$,
%OLD% $\lim_{Y\to\infty}\frac1Y\int_0^Y f(E_m(e^u))\,du
%OLD% =\int_{\mathbb R}f\,d\nu_m$. Under LI$(\chi_4)$, $\nu_m$ is the law of
%OLD% $X_m$.
%OLD% \item[(ii)] Under LI$(\chi_4)$, $\nu_m$ is absolutely continuous with
%OLD% support $\mathbb R$; the logarithmic density
%OLD% $\delta(m)$ of the set $\{x\ge2: R_m(x)>0\}$ exists, equals
%OLD% $\nu_m\bigl((0,\infty)\bigr)$, satisfies
%OLD% $\tfrac12<\delta(m)<1$ for every $m\in(-\tfrac12,\infty)$, and
%OLD% $m\mapsto\delta(m)$ is continuous on $(-\tfrac12,\infty)$.
%OLD% \item[(iii)] Under LI$(\chi_4)$, for all $m\in(-\tfrac12,\infty)$,
%OLD% \[
%OLD% 1-\delta(m)\;\le\;C\,(2m+1)^2
%OLD% \qquad\text{and}\qquad
%OLD% 1-\delta(m)\;\le\;\exp\!\Bigl(-\frac{1}{2C\,(2m+1)^2}\Bigr),
%OLD% \]
%OLD% with no restriction on $m$; the second bound is in fact the stronger of
%OLD% the two for every $m$. In particular $\delta(m)\to1$ as
%OLD% $m\to-\tfrac12^{\,+}$. Both bounds are informative only near the left
%OLD% endpoint: the (stronger) exponential bound drops below the trivial
%OLD% estimate $1-\delta(m)<\tfrac12$ of part (ii) precisely for
%OLD% $m<m_0:=\tfrac12\bigl((2C\log2)^{-1/2}-1\bigr)=0.57507\ldots$, and the
%OLD% first bound is $\ge1$ for $m\ge\tfrac12(C^{-1/2}-1)=0.76580\ldots$;
%OLD% for $m\ge m_0$ part (iii) asserts nothing beyond part (ii), and this
%OLD% theorem makes no claim about the rate at which $\delta(m)$ approaches
%OLD% $\tfrac12$ as $m\to\infty$.
%OLD% \item[(iv)] Consequently the classical density formulation of
%OLD% Chebyshev's bias at $m=0$ (logarithmic density
%OLD% $\delta(0)=0.995928\ldots$; $0.9959$ in \cite{RS94}, six decimals as in
%OLD% \cite{FM13} and re-derived here) is joined by the continuous family
%OLD% $\{\delta(m)\}_{m\in(-1/2,\infty)}$, with $\delta(m)\to1$ as
%OLD% $m\to-\tfrac12^{\,+}$, to the
%OLD% total-bias regime at the weighted threshold $m=-\tfrac12$, where the
%OLD% unlogged race $D_{-1/2}(x)=\sum_{p\le x}(-\chi_4(p))p^{-1/2}$ equals
%OLD% $\tfrac12\log\log x+O(1)$ and is eventually of one sign (under DRH
%OLD% \cite{AK23}; under GRH off a set of finite logarithmic measure
%OLD% \cite{Sheth24}; single-signed on a set of logarithmic density one under
%OLD% GRH alone \cite{Hay25}). None of these three endpoint results assumes
%OLD% LI. In the opposite direction the family now covers every positive
%OLD% weight; there the computed densities of Table~\ref{tab:delta} decrease
%OLD% toward $\tfrac12$, in the manner of the dissolution law of the Results
%OLD% section, which this theorem does not prove.
%OLD% \end{enumerate}
%OLD% \end{theorem}
%---NEW--->
\begin{theorem}\label{thm:interp}
Let $\chi$ be a real primitive character mod $q$. Assume GRH$(\chi)$
and NCZ$(\chi)$: $L(\tfrac12,\chi)\ne0$ (the latter unconditional for
$q=3,4$, as noted above).
\begin{enumerate}
\item[(i)] For each $m\in(-\tfrac12,\infty)$ the function
$u\mapsto E_m(e^u)$ possesses a limiting distribution
$\nu_m$: for every bounded continuous $f:\mathbb R\to\mathbb R$,
$\lim_{Y\to\infty}\frac1Y\int_0^Y f(E_m(e^u))\,du
=\int_{\mathbb R}f\,d\nu_m$. Under LI$(\chi)$, $\nu_m$ is the law of
$X_m$.
\item[(ii)] Under LI$(\chi)$, $\nu_m$ is absolutely continuous with
support $\mathbb R$; the logarithmic density
$\delta_\chi(m)$ of the set $\{x\ge2: R_m(x)>0\}$ exists, equals
$\nu_m\bigl((0,\infty)\bigr)$, satisfies
$\tfrac12<\delta_\chi(m)<1$ for every $m\in(-\tfrac12,\infty)$, and
$m\mapsto\delta_\chi(m)$ is continuous on $(-\tfrac12,\infty)$. (We
write $\delta(m)$ for $\delta_\chi(m)$ when $\chi$ is understood; for
$\chi=\chi_4,\chi_3$ these are the densities printed as
$\delta_4,\delta_3$ in Table~\ref{tab:delta}, and under LI$(\chi)$
they coincide with the quantities
$\delta_q(m)=\mathbb P(X_{m,\chi}>0)$ of
Section~\ref{sec:dissolution}.)
\item[(iii)] Under LI$(\chi)$, for all $m\in(-\tfrac12,\infty)$,
\[
1-\delta_\chi(m)\;\le\;C(\chi)\,(2m+1)^2
\qquad\text{and}\qquad
1-\delta_\chi(m)\;\le\;\exp\!\Bigl(-\frac{1}{2C(\chi)\,(2m+1)^2}\Bigr),
\]
with no restriction on $m$; the second bound is in fact the stronger of
the two for every $m$. In particular $\delta_\chi(m)\to1$ as
$m\to-\tfrac12^{\,+}$. Both bounds are informative only near the left
endpoint: the (stronger) exponential bound drops below the trivial
estimate $1-\delta_\chi(m)<\tfrac12$ of part (ii) precisely for
$m<m_0(\chi):=\tfrac12\bigl((2C(\chi)\log2)^{-1/2}-1\bigr)$, and the
first bound is $\ge1$ for $m\ge\tfrac12\bigl(C(\chi)^{-1/2}-1\bigr)$;
for $m\ge m_0(\chi)$ part (iii) asserts nothing beyond part (ii), and
this theorem makes no claim about the rate at which $\delta_\chi(m)$
approaches $\tfrac12$ as $m\to\infty$ (for that rate, see
Section~\ref{sec:dissolution}).
\item[(iv)] ($\chi=\chi_4$.) Consequently the classical density
formulation of
Chebyshev's bias at $m=0$ (logarithmic density
$\delta(0)=0.995928\ldots$; $0.9959$ in \cite{RS94}, six decimals as in
\cite{FM13} and re-derived here) is joined by the continuous family
$\{\delta(m)\}_{m\in(-1/2,\infty)}$, with $\delta(m)\to1$ as
$m\to-\tfrac12^{\,+}$, to the
total-bias regime at the weighted threshold $m=-\tfrac12$, where the
unlogged race $D_{-1/2}(x)=\sum_{p\le x}(-\chi_4(p))p^{-1/2}$ equals
$\tfrac12\log\log x+O(1)$ and is eventually of one sign (under DRH
\cite{AK23}; under GRH off a set of finite logarithmic measure
\cite{Sheth24}; single-signed on a set of logarithmic density one under
GRH alone \cite{Hay25}). None of these three endpoint results assumes
LI. In the opposite direction the family now covers every positive
weight; there the computed densities of Table~\ref{tab:delta} decrease
toward $\tfrac12$, in the manner of the dissolution law of the Results
section, which this theorem does not prove (it is proved, with a rate,
in Section~\ref{sec:dissolution}).
\end{enumerate}
\end{theorem}

\medskip\noindent\emph{Example ($q=4$).} $\chi_4$ is odd,
NCZ$(\chi_4)$ is the alternating-series positivity
$L(\tfrac12,\chi_4)>0$, and $C(\chi_4)=0.15602964988964488\ldots$ (in
closed form above). The thresholds of (iii) are
$m_0(\chi_4)=0.57507\ldots$ and
$\tfrac12\bigl(C(\chi_4)^{-1/2}-1\bigr)=0.76580\ldots$, the Chebyshev
bound falling below $\tfrac12$ only for $m<0.395058\ldots$; parts (ii)
and (iv) recover the classical anchor $\delta(0)=0.995928\ldots$ of
\cite{RS94,FM13} at $m=0$ (Table~\ref{tab:delta}).

\medskip\noindent\emph{Example ($q=3$).} The race $R_{m,\chi_3}$ pits
$p\equiv2$ against $p\equiv1\ (\mathrm{mod}\ 3)$; $\chi_3$ is odd,
NCZ$(\chi_3)$ is the grouped-series positivity
$L(\tfrac12,\chi_3)>0$, and
$C(\chi_3)=0.11340215685275857\ldots$ (Example above). Under
GRH$(\chi_3)$ and LI$(\chi_3)$, parts (i)--(iii) give existence,
continuity and the strict bounds for $\delta_3(m)$ on all of
$(-\tfrac12,\infty)$; the computed anchors are $\delta_3(0)=0.999063$
(the classical mod-$3$ bias, re-derived to six decimals in
Table~\ref{tab:delta}) and, from the sieves of the introduction, the
weighted first flips at $p=13$ ($m=1$) and $p=7$ ($m=2,3$), the
classical race $m=0$ not flipping below $10^{10}$.

\begin{corollary}[mod-$3$ rate]\label{cor:mod3rate}
Assume GRH$(\chi_3)$ and LI$(\chi_3)$ (NCZ$(\chi_3)$ holds
unconditionally). Then the logarithmic density
$\delta_3(m)=\delta_{\chi_3}(m)$ of the set
$\{x\ge2:R_{m,\chi_3}(x)>0\}$ satisfies, for every
$m\in(-\tfrac12,\infty)$,
\[
1-\delta_3(m)\;\le\;
\exp\!\Bigl(-\frac{1}{2\,C(\chi_3)\,(2m+1)^2}\Bigr),
\qquad
C(\chi_3)=0.11340215685275857\ldots,
\]
and likewise $1-\delta_3(m)\le C(\chi_3)(2m+1)^2$; in particular
$\delta_3(m)\to1$ as $m\to-\tfrac12^{\,+}$. The exponential bound is
informative (below $\tfrac12$) precisely for
$m<m_0(\chi_3)=0.761048\ldots$, the Chebyshev form for
$m<0.549891\ldots$; at $m=0$ the exponential bound gives
$1-\delta_3(0)\le e^{-1/(2C(\chi_3))}=0.01216\ldots$, against the
computed value $9.37\times10^{-4}$ (Table~\ref{tab:delta}).
\end{corollary}

\begin{proof}
Theorem~\ref{thm:interp}(iii) with $\chi=\chi_3$; the numerical
thresholds are evaluated from $C(\chi_3)$ as in the Example above.
\end{proof}
%<---NEW---

% ============================================================================
% Remark after the theorem (Humphries analogy): C -> C(chi).
% Lines 424-425.
% ============================================================================
%% BLOCK 5 REPLACES:
%OLD% normalized mean is fixed at $1$ while the variance collapses like
%OLD% $C(2m+1)^2$; in both cases the bias-to-spread ratio diverges, and
%---NEW--->
normalized mean is fixed at $1$ while the variance collapses like
$C(\chi)(2m+1)^2$; in both cases the bias-to-spread ratio diverges, and
%<---NEW---

% ============================================================================
% Lemma lem:classical restated for a real primitive chi of either
% parity; for even chi the s=0 trivial zero contributes -log x + O_q(1),
% absorbed by the existing floor.  Lines 435-445.
% ============================================================================
%% BLOCK 6 REPLACES:
%OLD% \begin{lemma}[truncated explicit formula, all truncations]\label{lem:classical}
%OLD% Let $\chi=\chi_4$ and let $\rho=\beta+i\gamma$ run over the nontrivial
%OLD% zeros of $L(s,\chi)$, listed with multiplicity. Unconditionally, for all
%OLD% $x\ge2$ and all $T\ge2$,
%OLD% \begin{equation}\label{eq:classical}
%OLD% \psi(x,\chi)\;=\;-\sum_{|\gamma|\le T}\frac{x^{\rho}}{\rho}
%OLD% +O\!\Bigl(\frac{x\log^2(xT)}{T}+\log x\Bigr).
%OLD% \end{equation}
%OLD% No inequality between $T$ and $x$ is assumed; in particular $T>x$ is
%OLD% allowed.
%OLD% \end{lemma}
%---NEW--->
\begin{lemma}[truncated explicit formula, all truncations]\label{lem:classical}
Let $\chi$ be a real primitive character mod $q$, of parity
$\chi(-1)=(-1)^a$, and let $\rho=\beta+i\gamma$ run over the nontrivial
zeros of $L(s,\chi)$, listed with multiplicity. Unconditionally, for all
$x\ge2$ and all $T\ge2$,
\begin{equation}\label{eq:classical}
\psi(x,\chi)\;=\;-\sum_{|\gamma|\le T}\frac{x^{\rho}}{\rho}
+O_q\!\Bigl(\frac{x\log^2(xT)}{T}+\log x\Bigr).
\end{equation}
No inequality between $T$ and $x$ is assumed; in particular $T>x$ is
allowed. For even $\chi$ ($a=0$) the trivial zero of $L(s,\chi)$ at
$s=0$ contributes a genuine $-\log x+O_q(1)$ term, absorbed by the
$O_q(\log x)$ floor (cf.\ \cite[Ch.~19]{Dav00}); this is the only
point at which the parity alters the shape of the formula.
\end{lemma}
%<---NEW---

% ============================================================================
% Proof of lem:classical, opening: constants now depend on q.  Line 448.
% ============================================================================
%% BLOCK 7 REPLACES:
%OLD% All implied constants are absolute ($q=4$ is fixed). We use five
%---NEW--->
All implied constants depend at most on $q$ ($\chi$ is fixed). We use
five
%<---NEW---

% ============================================================================
% (F3) for general chi.  Lines 482-485.
% ============================================================================
%% BLOCK 8 REPLACES:
%OLD% \emph{(F3)} $N(t+1,\chi)-N(t,\chi)\ll\log(t+2)$ for $t\ge0$, where
%OLD% $N(t,\chi)$ counts zeros with $0<\gamma\le t$; unconditional
%OLD% Riemann--von Mangoldt for $L(s,\chi_4)$ \cite[\S16]{Dav00}. Zeros are
%OLD% symmetric under $\gamma\mapsto-\gamma$ ($\chi_4$ is real).
%---NEW--->
\emph{(F3)} $N(t+1,\chi)-N(t,\chi)\ll_q\log(t+2)$ for $t\ge0$, where
$N(t,\chi)$ counts zeros with $0<\gamma\le t$; unconditional
Riemann--von Mangoldt for $L(s,\chi)$ \cite[\S16]{Dav00}. Zeros are
symmetric under $\gamma\mapsto-\gamma$ ($\chi$ is real).
%<---NEW---

% ============================================================================
% (F5) and the s=0 analysis for both parities: functional equation with
% the parity-dependent Gamma factor; odd chi has L(0,chi)!=0, even chi
% has a simple trivial zero at s=0 with L'/L(s)=1/s+b(chi)+O(s).
% Lines 492-521.
% ============================================================================
%% BLOCK 9 REPLACES:
%OLD% \emph{(F5)} $\Lambda(s)=(4/\pi)^{(s+1)/2}\Gamma(\tfrac{s+1}2)L(s,\chi)$
%OLD% satisfies $\Lambda(s)=\Lambda(1-s)$ (\cite[\S9]{Dav00}; the root number
%OLD% of $\chi_4$ is $\tau(\chi_4)/(i\sqrt4)=2i/2i=+1$). Taking logarithmic
%OLD% derivatives,
%OLD% \[
%OLD% \frac{L'}{L}(s,\chi)
%OLD% =-\log\frac4\pi-\tfrac12\psi\Bigl(\frac{s+1}2\Bigr)
%OLD% -\tfrac12\psi\Bigl(\frac{2-s}2\Bigr)-\frac{L'}{L}(1-s,\chi).
%OLD% \]
%OLD% If $\sigma\le-1$ and either $|t|\ge2$ or $\sigma$ is a negative even
%OLD% integer, then: $\mathrm{Re}(1-s)\ge2$, so $|L'/L(1-s,\chi)|=O(1)$ by
%OLD% (F4); $\mathrm{Re}\bigl(\frac{2-s}2\bigr)\ge\tfrac32$, so
%OLD% $|\psi(\frac{2-s}2)|\ll\log(2+|s|)$ by Stirling; and for
%OLD% $w=\frac{s+1}2$ the reflection $\psi(w)=\psi(1-w)-\pi\cot(\pi w)$
%OLD% applies with $\mathrm{Re}(1-w)\ge1$ and with $|\cot(\pi w)|\le
%OLD% \coth(\pi)$ when $|\mathrm{Im}\,w|=|t|/2\ge1$, respectively
%OLD% $|\cot(\pi w)|=|\tanh(\pi t/2)|\le1$ when $\mathrm{Re}\,w$ is a
%OLD% half-odd-integer (which is the case when $\sigma$ is a negative even
%OLD% integer). Hence on those sets
%OLD% $|L'/L(s,\chi)|\ll\log(2+|s|)$.
%OLD%
%OLD% Also $L(0,\chi_4)=\tfrac12\ne0$: evaluating $\Lambda(0)=\Lambda(1)$
%OLD% from (F5),
%OLD% $(4/\pi)^{1/2}\Gamma(\tfrac12)L(0,\chi_4)=(4/\pi)\Gamma(1)L(1,\chi_4)$,
%OLD% and $L(1,\chi_4)=\pi/4$ (Leibniz), so
%OLD% $L(0,\chi_4)=(4/\pi)^{1/2}(\pi/4)/\sqrt{\pi}=\tfrac12$.
%OLD% Consequently $L'/L$ is regular at $s=0$. Moreover every nontrivial zero
%OLD% has $0<\beta<1$: there are no zeros with $\beta=1$ (de la Vall\'ee
%OLD% Poussin, unconditional), hence none with $\beta=0$ by (F5); so the pole
%OLD% of $1/s$ at $s=0$ collides with no zero.
%---NEW--->
\emph{(F5)} $\Lambda(s,\chi)=(q/\pi)^{(s+a)/2}
\Gamma\bigl(\tfrac{s+a}2\bigr)L(s,\chi)$
satisfies $\Lambda(s,\chi)=\Lambda(1-s,\chi)$: the root number of a
real primitive character is $+1$ (\cite[Chs.~2, 9]{Dav00}; Gauss's
evaluation $\tau(\chi)=\sqrt q$ for $a=0$, $i\sqrt q$ for $a=1$ ---
for $q=4$ the two-term sum gives $\tau(\chi_4)/(i\sqrt4)=2i/2i=+1$).
Taking logarithmic derivatives,
\[
\frac{L'}{L}(s,\chi)
=-\log\frac q\pi-\tfrac12\psi\Bigl(\frac{s+a}2\Bigr)
-\tfrac12\psi\Bigl(\frac{1-s+a}2\Bigr)-\frac{L'}{L}(1-s,\chi).
\]
If $\sigma\le-1$ and either $|t|\ge2$ or $\sigma+a$ is an odd
(negative) integer, then: $\mathrm{Re}(1-s)\ge2$, so
$|L'/L(1-s,\chi)|=O(1)$ by
(F4); $\mathrm{Re}\bigl(\frac{1-s+a}2\bigr)\ge1$, so
$|\psi(\frac{1-s+a}2)|\ll\log(2+|s|)$ by Stirling; and for
$w=\frac{s+a}2$ the reflection $\psi(w)=\psi(1-w)-\pi\cot(\pi w)$
applies with $\mathrm{Re}(1-w)\ge1$ and with $|\cot(\pi w)|\le
\coth(\pi)$ when $|\mathrm{Im}\,w|=|t|/2\ge1$, respectively
$|\cot(\pi w)|=|\tanh(\pi t/2)|\le1$ when $\mathrm{Re}\,w$ is a
half-odd-integer (which is the case when $\sigma+a$ is an odd
integer). Hence on those sets
$|L'/L(s,\chi)|\ll_q\log(2+|s|)$.

The behaviour at $s=0$ depends on the parity. For odd $\chi$ ($a=1$),
$L(0,\chi)\ne0$: evaluating $\Lambda(0,\chi)=\Lambda(1,\chi)$
from (F5),
$(q/\pi)^{1/2}\Gamma(\tfrac12)L(0,\chi)=(q/\pi)\Gamma(1)L(1,\chi)$
with $L(1,\chi)\ne0$ (Dirichlet); for $\chi_4$, $L(1,\chi_4)=\pi/4$
(Leibniz) gives the explicit value
$L(0,\chi_4)=(4/\pi)^{1/2}(\pi/4)/\sqrt{\pi}=\tfrac12$.
Consequently, for odd $\chi$, $L'/L$ is regular at $s=0$. For even
$\chi$ ($a=0$), $s=0$ is a \emph{simple} trivial zero of $L(s,\chi)$:
$\Gamma(s/2)$ has a simple pole there while
$\Lambda(0,\chi)=\Lambda(1,\chi)\ne0$ is finite, so
\[
\frac{L'}{L}(s,\chi)\;=\;\frac1s+b(\chi)+O(|s|)\qquad(s\to0)
\]
for a finite constant $b(\chi)=O_\chi(1)$ (cf.\ \cite[Ch.~19]{Dav00}).
Moreover every nontrivial zero
has $0<\beta<1$: there are no zeros with $\beta=1$ (Dirichlet, de la
Vall\'ee Poussin, unconditional), hence, by (F5), none with $\beta=0$
apart from the trivial zero at $s=0$ in the even case, which is
accounted for separately below.
%<---NEW---

% ============================================================================
% Step 3 of the contour argument, both parities: U=2K+1-a so the left
% edge passes between trivial zeros; for even chi the double pole at
% s=0 has residue -log x - b(chi).  Lines 564-577.
% ============================================================================
%% BLOCK 10 REPLACES:
%OLD% \emph{Step 3 (contour).} For $K\in\mathbb N$ set $U=2K$ and integrate
%OLD% $(-L'/L)(s,\chi)\,x^s/s$ counterclockwise around the rectangle with
%OLD% vertices $c\pm iT'$, $-U\pm iT'$. The poles inside are: the nontrivial
%OLD% zeros with $|\gamma|\le T'$ (residue $-x^\rho/\rho$ each, with
%OLD% multiplicity; $\rho\ne0$ since $0<\beta<1$); the simple pole of $1/s$
%OLD% at $s=0$ (residue $-L'/L(0,\chi)=O(1)$, using $L(0,\chi)=\tfrac12$);
%OLD% and the trivial zeros $s=-(2j+1)$, $0\le j\le K-1$, of $L(s,\chi_4)$
%OLD% (simple, from the poles of $\Gamma(\tfrac{s+1}2)$; residue
%OLD% $+x^{-(2j+1)}/(2j+1)$ each). Hence
%OLD% \[
%OLD% \frac{1}{2\pi i}\oint
%OLD% =-\sum_{|\gamma|\le T'}\frac{x^\rho}{\rho}-\frac{L'}{L}(0,\chi)
%OLD% +\sum_{j=0}^{K-1}\frac{x^{-(2j+1)}}{2j+1}.
%OLD% \]
%---NEW--->
\emph{Step 3 (contour).} For $K\in\mathbb N$ set $U=2K+1-a$ (so the
line $\sigma=-U$ passes midway between consecutive trivial zeros in
either parity) and integrate
$(-L'/L)(s,\chi)\,x^s/s$ counterclockwise around the rectangle with
vertices $c\pm iT'$, $-U\pm iT'$. The poles inside are: the nontrivial
zeros with $|\gamma|\le T'$ (residue $-x^\rho/\rho$ each, with
multiplicity; $\rho\ne0$ since $0<\beta<1$); the point $s=0$; and the
trivial zeros $s=-(2j+a)$ of $L(s,\chi)$ with $j\ge1-a$ and
$2j+a\le U$ (simple, from the poles of $\Gamma(\tfrac{s+a}2)$; residue
$+x^{-(2j+a)}/(2j+a)$ each). At $s=0$: for odd $\chi$ the integrand
has only the simple pole of $1/s$, with residue
$-L'/L(0,\chi)=O_q(1)$ (using $L(0,\chi)\ne0$); for even $\chi$ the
pole of $1/s$ meets the simple trivial zero of $L$, and the Laurent
expansion of (F5) gives a double pole with residue
\[
\mathrm{Res}_{s=0}\,\Bigl(-\frac{L'}{L}(s,\chi)\Bigr)\frac{x^s}{s}
\;=\;-\log x-b(\chi)\;=\;O_q(\log x).
\]
Writing $r_0(x,\chi)$ for the residue at $s=0$ just computed (in
either parity $r_0(x,\chi)=O_q(\log x)$),
\[
\frac{1}{2\pi i}\oint
=-\sum_{|\gamma|\le T'}\frac{x^\rho}{\rho}+r_0(x,\chi)
+\sum_{\substack{j\ge1-a\\ 2j+a\le U}}\frac{x^{-(2j+a)}}{2j+a}.
\]
%<---NEW---

% ============================================================================
% Step 6, both parities: the limiting trivial-zero sums are O(1) in
% either parity, and the O_q(log x) floor carries the even-chi
% -log x - b(chi) term.  Lines 603-610.
% ============================================================================
%% BLOCK 11 REPLACES:
%OLD% \emph{Step 6 (assemble; $U\to\infty$).} Letting $K\to\infty$, the
%OLD% trivial-zero sum increases to
%OLD% $\sum_{j\ge0}x^{-(2j+1)}/(2j+1)=\operatorname{artanh}(1/x)
%OLD% \le\operatorname{artanh}(\tfrac12)<0.55$. Combining Steps 1--5,
%OLD% \[
%OLD% \psi(x,\chi)=-\sum_{|\gamma|\le T'}\frac{x^\rho}{\rho}
%OLD% +O\Bigl(\frac{x\log^2(xT)}{T}+\log x\Bigr).
%OLD% \]
%---NEW--->
\emph{Step 6 (assemble; $U\to\infty$).} Letting $K\to\infty$, the
trivial-zero sum increases to
$\sum_{j\ge0}x^{-(2j+1)}/(2j+1)=\operatorname{artanh}(1/x)
\le\operatorname{artanh}(\tfrac12)<0.55$ for odd $\chi$, and to
$\sum_{j\ge1}x^{-2j}/(2j)=-\tfrac12\log(1-x^{-2})
\le\tfrac12\log\tfrac43<0.15$ for even $\chi$ (both for $x\ge2$);
and $r_0(x,\chi)=O_q(\log x)$ in either parity --- for even $\chi$
the floor $O_q(\log x)$ is genuinely attained, by the $-\log x$ of
Step 3. Combining Steps 1--5,
\[
\psi(x,\chi)=-\sum_{|\gamma|\le T'}\frac{x^\rho}{\rho}
+O_q\Bigl(\frac{x\log^2(xT)}{T}+\log x\Bigr).
\]
%<---NEW---

% ============================================================================
% Lemma lem:trunc restated under GRH(chi)+NCZ(chi).  Lines 620-631.
% ============================================================================
%% BLOCK 12 REPLACES:
%OLD% \begin{lemma}[shifted truncated explicit formula]\label{lem:trunc}
%OLD% Assume GRH$(\chi_4)$ and let $m\in(-\tfrac12,\infty)$ be fixed. For
%OLD% $\psi_m(x):=\sum_{n\le x}\Lambda(n)\chi(n)n^m$ and all $x\ge2$,
%OLD% $T\ge2$ (no inequality between $T$ and $x$ is assumed;
%OLD% Lemma~\ref{lem:classical}),
%OLD% \begin{equation}\label{eq:trunc}
%OLD% \psi_m(x)\;=\;-\sum_{|\gamma|\le T}\frac{x^{\rho+m}}{\rho+m}
%OLD% \;+\;O_m\!\Bigl(\frac{x^{m+1}\log^2(xT)}{T}+(1+x^m)\log x\Bigr).
%OLD% \end{equation}
%OLD% For $m\le0$ the floor $(1+x^m)\log x\le2\log x$ recovers the earlier
%OLD% form of this lemma.
%OLD% \end{lemma}
%---NEW--->
\begin{lemma}[shifted truncated explicit formula]\label{lem:trunc}
Assume GRH$(\chi)$ and NCZ$(\chi)$, and let $m\in(-\tfrac12,\infty)$
be fixed. For
$\psi_m(x):=\sum_{n\le x}\Lambda(n)\chi(n)n^m$ and all $x\ge2$,
$T\ge2$ (no inequality between $T$ and $x$ is assumed;
Lemma~\ref{lem:classical}),
\begin{equation}\label{eq:trunc}
\psi_m(x)\;=\;-\sum_{|\gamma|\le T}\frac{x^{\rho+m}}{\rho+m}
\;+\;O_{m,\chi}\!\Bigl(\frac{x^{m+1}\log^2(xT)}{T}+(1+x^m)\log x\Bigr).
\end{equation}
For $m\le0$ the floor $(1+x^m)\log x\le2\log x$ recovers the earlier
form of this lemma.
\end{lemma}
%<---NEW---

% ============================================================================
% Proof of lem:trunc: the numeric bound gamma_1 > 6 (true for chi_4)
% is replaced by gamma_1(chi) > 0, which is exactly GRH + NCZ; the
% constants now depend on chi.  Lines 648-657.
% ============================================================================
%% BLOCK 13 REPLACES:
%OLD% an identity (integrate $t^{\rho+m-1}$; note
%OLD% $\mathrm{Re}(\rho+m)=m+\tfrac12>0$ for every $m>-\tfrac12$, and
%OLD% $|\rho+m|\ge\gamma\ge\gamma_1>6$, so no denominator degenerates at any
%OLD% $m$ in the range). The endpoint terms are \emph{summable}: since
%OLD% $|\rho|\ge\gamma$ and $|\rho+m|\ge\gamma$,
%OLD% \[
%OLD% \sum_{|\gamma|\le T}\Bigl|\frac{m\,2^{\rho+m}}{\rho(\rho+m)}\Bigr|
%OLD% \;\le\;|m|\,2^{m+1/2}\sum_{|\gamma|\le T}\frac{1}{\gamma^{2}}
%OLD% \;\ll_m\;1,
%OLD% \]
%---NEW--->
an identity (integrate $t^{\rho+m-1}$; note
$\mathrm{Re}(\rho+m)=m+\tfrac12>0$ for every $m>-\tfrac12$, and
$|\rho+m|\ge|\gamma|\ge\gamma_1(\chi)>0$ --- every ordinate is nonzero
because under GRH a zero with $\gamma=0$ would sit at $s=\tfrac12$,
excluded by NCZ$(\chi)$ --- so no denominator degenerates at any
$m$ in the range; the constants below depend on $\chi$ through
$\gamma_1(\chi)$ and $\sum_\gamma\gamma^{-2}$). The endpoint terms are
\emph{summable}: since
$|\rho|\ge|\gamma|$ and $|\rho+m|\ge|\gamma|$,
\[
\sum_{|\gamma|\le T}\Bigl|\frac{m\,2^{\rho+m}}{\rho(\rho+m)}\Bigr|
\;\le\;|m|\,2^{m+1/2}\sum_{|\gamma|\le T}\frac{1}{\gamma^{2}}
\;\ll_{m,\chi}\;1,
\]
%<---NEW---

% ============================================================================
% Lemma lem:squares statement: O_{q,m} error (the p | q prime-square
% terms are O_{q,m}(1) constants, absorbed).  Lines 702-707.
% ============================================================================
%% BLOCK 14 REPLACES:
%OLD% \begin{lemma}[prime-power separation]\label{lem:squares}
%OLD% Fix $m\in(-\tfrac12,\infty)$. For $x\ge4$,
%OLD% \[
%OLD% -R_m(x)\;=\;\psi_m(x)\;-\;\sum_{p\le\sqrt x}p^{2m}\log p
%OLD% \;+\;O_m\!\bigl(x^{m+1/3}\log x+\log^2x\bigr),
%OLD% \]
%---NEW--->
\begin{lemma}[prime-power separation]\label{lem:squares}
Fix $m\in(-\tfrac12,\infty)$. For $x\ge4$,
\[
-R_m(x)\;=\;\psi_m(x)\;-\;\sum_{p\le\sqrt x}p^{2m}\log p
\;+\;O_{q,m}\!\bigl(x^{m+1/3}\log x+\log^2x\bigr),
\]
%<---NEW---

% ============================================================================
% Proof of lem:squares, k=2 part: sum over p coprime to q; the at most
% omega(q) primes p | q each contribute p^{2m} log p <= (1+q^{2m}) log q
% = O_{q,m}(1) -- constants, NOT growing in x.  Lines 722-726.
% ============================================================================
%% BLOCK 15 REPLACES:
%OLD% $\psi_m$ splits over prime powers $p^k\le x$. The $k=1$ part is
%OLD% $\sum_{p\le x}\chi(p)p^m\log p=-R_m(x)$. The $k=2$ part is
%OLD% $\sum_{p^2\le x}\chi(p)^2p^{2m}\log p=\sum_{3\le p\le\sqrt x}p^{2m}\log p
%OLD% =\sum_{p\le\sqrt x}p^{2m}\log p+O_m(1)$, since $\chi(p)^2=1$ for odd
%OLD% $p$ and $\chi(2)=0$.
%---NEW--->
$\psi_m$ splits over prime powers $p^k\le x$. The $k=1$ part is
$\sum_{p\le x}\chi(p)p^m\log p=-R_m(x)$. The $k=2$ part is
\[
\sum_{p^2\le x}\chi(p)^2p^{2m}\log p
\;=\;\sum_{\substack{p\le\sqrt x\\ p\nmid q}}p^{2m}\log p
\;=\;\sum_{p\le\sqrt x}p^{2m}\log p+O_{q,m}(1),
\]
since $\chi(p)^2=1$ for $p\nmid q$ and $\chi(p)=0$ for $p\mid q$: the
subtracted terms are the at most $\omega(q)$ primes dividing $q$, each
contributing $p^{2m}\log p\le(1+q^{2m})\log q$, a constant depending
only on $q$ and $m$ (for $q=4$ this is the single term $p=2$, which is
absent from both sides).
%<---NEW---

% ============================================================================
% Proof of (i): conjugation symmetry for a general real chi.
% Lines 824-827.
% ============================================================================
%% BLOCK 16 REPLACES:
%OLD% Pairing $\rho$ and $\bar\rho$ in Lemma~\ref{lem:trunc} (the zeros of
%OLD% $L(s,\chi_4)$ are symmetric under conjugation since $\chi_4$ is real)
%OLD% and combining with Lemma~\ref{lem:squares}, for $u\ge2$ and any
%OLD% $T\ge2$,
%---NEW--->
Pairing $\rho$ and $\bar\rho$ in Lemma~\ref{lem:trunc} (the zeros of
$L(s,\chi)$ are symmetric under conjugation since $\chi$ is real)
and combining with Lemma~\ref{lem:squares}, for $u\ge2$ and any
$T\ge2$,
%<---NEW---

% ============================================================================
% Proof of (i): Riemann-von Mangoldt with the conductor q.
% Lines 863-868.
% ============================================================================
%% BLOCK 17 REPLACES:
%OLD% The coefficient conditions: $|r_\gamma(m)|\le2(2m+1)/\gamma$, so
%OLD% $\beta=1>\tfrac12$; and $\sum_{T<\gamma\le T+1}1\ll\log T$ is the
%OLD% Riemann--von Mangoldt formula for $L(s,\chi_4)$,
%OLD% $N(T,\chi_4)=\frac{T}{2\pi}\log\frac{2T}{\pi e}+O(\log T)$
%OLD% (\cite[\S16]{Dav00}), which holds unconditionally. Hence
%OLD% \cite[Cor.~1.3(a)]{ANS14} applies verbatim and $\nu_m$ exists. (The
%---NEW--->
The coefficient conditions: $|r_\gamma(m)|\le2(2m+1)/\gamma$, so
$\beta=1>\tfrac12$; and $\sum_{T<\gamma\le T+1}1\ll_q\log T$ is the
Riemann--von Mangoldt formula for $L(s,\chi)$,
$N(T,\chi)=\frac{T}{2\pi}\log\frac{qT}{2\pi e}+O_q(\log T)$
(\cite[\S16]{Dav00}), which holds unconditionally. Hence
\cite[Cor.~1.3(a)]{ANS14} applies verbatim and $\nu_m$ exists. (The
%<---NEW---

% ============================================================================
% Proof of (i): LI(chi_4) -> LI(chi).  Lines 876-878.
% ============================================================================
%% BLOCK 18 REPLACES:
%OLD% Under LI$(\chi_4)$ the multiset $\{\gamma\}$ is linearly independent
%OLD% over $\mathbb Q$, and \cite[Thm.~1.9]{ANS14} (proved in their \S5)
%OLD% gives
%---NEW--->
Under LI$(\chi)$ the multiset $\{\gamma\}$ is linearly independent
over $\mathbb Q$, and \cite[Thm.~1.9]{ANS14} (proved in their \S5)
gives
%<---NEW---

% ============================================================================
% Proof of (ii), fact (beta): zero-counting with the conductor q.
% Lines 919-924.
% ============================================================================
%% BLOCK 19 REPLACES:
%OLD% ($\beta$) $\sum_nc_n=\infty$: for every $\gamma\ge m+\tfrac12$ --- all
%OLD% but finitely many ordinates --- one has
%OLD% $(m+\tfrac12)^2\le\gamma^2$, hence
%OLD% $a_\gamma(m)\ge(2m+1)/(\gamma\sqrt2)$, and
%OLD% $\sum_{\gamma\le T}\gamma^{-1}\gg(\log T)^2$ by partial summation
%OLD% against $N(T,\chi_4)\sim\frac{T}{2\pi}\log\frac{2T}{\pi e}$;
%---NEW--->
($\beta$) $\sum_nc_n=\infty$: for every $\gamma\ge m+\tfrac12$ --- all
but finitely many ordinates --- one has
$(m+\tfrac12)^2\le\gamma^2$, hence
$a_\gamma(m)\ge(2m+1)/(\gamma\sqrt2)$, and
$\sum_{\gamma\le T}\gamma^{-1}\gg_q(\log T)^2$ by partial summation
against $N(T,\chi)\sim\frac{T}{2\pi}\log\frac{qT}{2\pi e}$;
%<---NEW---

% ============================================================================
% Proof of (iii): C -> C(chi) in the variance and Chebyshev bounds.
% Lines 1013-1022.
% ============================================================================
%% BLOCK 20 REPLACES:
%OLD% \[
%OLD% \sigma_m^2:=\mathrm{Var}(V_m)=\sum_{\gamma>0}2a_\gamma(m)^2
%OLD% \;\le\;(2m+1)^2\sum_{\gamma>0}\frac{2}{\gamma^2}\;=\;C(2m+1)^2 ,
%OLD% \]
%OLD% since $a_\gamma(m)^2\le(2m+1)^2/\gamma^2$. Chebyshev's inequality
%OLD% gives
%OLD% \[
%OLD% 1-\delta(m)=\mathbb P(V_m\le-1)\le\mathbb P(|V_m|\ge1)\le\sigma_m^2
%OLD% \le C(2m+1)^2 .
%OLD% \]
%---NEW--->
\[
\sigma_m^2:=\mathrm{Var}(V_m)=\sum_{\gamma>0}2a_\gamma(m)^2
\;\le\;(2m+1)^2\sum_{\gamma>0}\frac{2}{\gamma^2}\;=\;C(\chi)(2m+1)^2 ,
\]
since $a_\gamma(m)^2\le(2m+1)^2/\gamma^2$. Chebyshev's inequality
gives
\[
1-\delta(m)=\mathbb P(V_m\le-1)\le\mathbb P(|V_m|\ge1)\le\sigma_m^2
\le C(\chi)(2m+1)^2 .
\]
%<---NEW---

% ============================================================================
% Proof of (iii): C -> C(chi) in the definition of A.  Lines 1026-1027.
% ============================================================================
%% BLOCK 21 REPLACES:
%OLD% For the exponential bound, set $A:=\sum_{\gamma>0}a_\gamma(m)^2
%OLD% =\sigma_m^2/2\le\tfrac{C}{2}(2m+1)^2$ and let
%---NEW--->
For the exponential bound, set $A:=\sum_{\gamma>0}a_\gamma(m)^2
=\sigma_m^2/2\le\tfrac{C(\chi)}{2}(2m+1)^2$ and let
%<---NEW---

% ============================================================================
% Proof of (iii): C -> C(chi) in the Chernoff conclusion.
% Lines 1049-1052.
% ============================================================================
%% BLOCK 22 REPLACES:
%OLD% \[
%OLD% 1-\delta(m)=\mathbb P(V_m\le-1)\;\le\;\exp\Bigl(-\frac1{4A}\Bigr)
%OLD% \;\le\;\exp\Bigl(-\frac{1}{2C(2m+1)^2}\Bigr).
%OLD% \]
%---NEW--->
\[
1-\delta(m)=\mathbb P(V_m\le-1)\;\le\;\exp\Bigl(-\frac1{4A}\Bigr)
\;\le\;\exp\Bigl(-\frac{1}{2C(\chi)(2m+1)^2}\Bigr).
\]
%<---NEW---

% ============================================================================
% Scope paragraph of the proof of (iii): thresholds as functions of
% C(chi) with the chi_4 and chi_3 numerics; and the STALE SENTENCE FIX
% (task (e)): the dissolution law IS proved in this paper
% (Section sec:dissolution / Theorem thm:dissolution).  Lines 1059-1078.
% ============================================================================
%% BLOCK 23 REPLACES:
%OLD% \emph{Scope.} Although valid on the whole range, the bounds carry
%OLD% information only near the left endpoint, and we state this openly.
%OLD% With $C=0.1560296\ldots$: the Chebyshev bound $C(2m+1)^2$ is below the
%OLD% trivial estimate $\tfrac12$ of part (ii) only for
%OLD% $m<\tfrac12\bigl((2C)^{-1/2}-1\bigr)=0.39506\ldots$ and is $\ge1$ for
%OLD% $m\ge\tfrac12\bigl(C^{-1/2}-1\bigr)=0.76580\ldots$; the exponential
%OLD% bound is below $\tfrac12$ only for
%OLD% $m<m_0=\tfrac12\bigl((2C\log2)^{-1/2}-1\bigr)=0.57507\ldots$. For
%OLD% $m\ge m_0$, part (iii) asserts nothing beyond part (ii). This
%OLD% limitation is consistent with the expected truth: the computed
%OLD% densities approach $\tfrac12$ from above (Table~\ref{tab:delta}), so an
%OLD% upper bound on $1-\delta(m)$ tending to $0$ for large $m$ would
%OLD% contradict the numerics, and none should be sought. The conjectured
%OLD% large-$m$ statement,
%OLD% $\delta(m)-\tfrac12\sim(2\pi\,\mathrm{Var}_4(m))^{-1/2}$ with
%OLD% $\mathrm{Var}_4(m)$ as in Results item 3 (asymptotically equal to
%OLD% $\sigma_m^2=\mathrm{Var}(V_m)$ above), is the
%OLD% dissolution law of the Results section; its proof is the subject of
%OLD% Theorem~\ref{thm:dissolution} of the companion draft and is \emph{not}
%OLD% proved in this paper.
%---NEW--->
\emph{Scope.} Although valid on the whole range, the bounds carry
information only near the left endpoint, and we state this openly.
The Chebyshev bound $C(\chi)(2m+1)^2$ is below the
trivial estimate $\tfrac12$ of part (ii) only for
$m<\tfrac12\bigl((2C(\chi))^{-1/2}-1\bigr)$ and is $\ge1$ for
$m\ge\tfrac12\bigl(C(\chi)^{-1/2}-1\bigr)$; the exponential
bound is below $\tfrac12$ only for
$m<m_0(\chi)=\tfrac12\bigl((2C(\chi)\log2)^{-1/2}-1\bigr)$. At the two
computed moduli: for $\chi_4$ ($C=0.1560296\ldots$) these thresholds
are $0.395058\ldots$, $0.76580\ldots$ and $m_0(\chi_4)=0.57507\ldots$;
for $\chi_3$ ($C=0.1134021\ldots$) they are $0.549891\ldots$,
$0.984770\ldots$ and $m_0(\chi_3)=0.761048\ldots$ For
$m\ge m_0(\chi)$, part (iii) asserts nothing beyond part (ii). This
limitation is consistent with the expected truth: the computed
densities approach $\tfrac12$ from above (Table~\ref{tab:delta}), so an
upper bound on $1-\delta(m)$ tending to $0$ for large $m$ would
contradict the numerics, and none should be sought. The large-$m$
statement,
$\delta_q(m)-\tfrac12\sim(2\pi\,\mathrm{Var}_q(m))^{-1/2}$ with
$\mathrm{Var}_q(m)$ as in Results item 3 (asymptotically equal to
$\sigma_m^2=\mathrm{Var}(V_m)$ above), is the
dissolution law of the Results section; it is not proved by this
theorem, but it \emph{is} proved in this paper:
Theorem~\ref{thm:dissolution} of Section~\ref{sec:dissolution}
establishes it, with rate and a second-order term, for the limiting
variable of every real primitive $\chi$ under GRH$(\chi)$ and
NCZ$(\chi)$, and, under LI$(\chi)$, part (ii) of the present theorem
transfers it to the logarithmic density of the race itself
(Corollary~\ref{cor:dissolution-race} records the case $q=4$).
%<---NEW---

% ============================================================================
% Remark on constants: C -> C(chi).  Lines 1080-1084.
% ============================================================================
%% BLOCK 24 REPLACES:
%OLD% \emph{Remark on constants.} Plain Hoeffding (range form:
%OLD% $\mathbb P(S\le-t)\le\exp(-2t^2/\sum_i(\text{range}_i)^2)$ with
%OLD% $\mathrm{range}_\gamma=4a_\gamma$) gives
%OLD% $\exp(-1/(8A))\le\exp\bigl(-1/(4C(2m+1)^2)\bigr)$, valid for all $m$ as
%OLD% well; the arcsine moment-generating function is exactly $I_0$, and
%---NEW--->
\emph{Remark on constants.} Plain Hoeffding (range form:
$\mathbb P(S\le-t)\le\exp(-2t^2/\sum_i(\text{range}_i)^2)$ with
$\mathrm{range}_\gamma=4a_\gamma$) gives
$\exp(-1/(8A))\le\exp\bigl(-1/(4C(\chi)(2m+1)^2)\bigr)$, valid for all
$m$ as
well; the arcsine moment-generating function is exactly $I_0$, and
%<---NEW---

% ============================================================================
% Numerical-consistency paragraph: scope it to chi = chi_4 (its numbers
% are the mod-4 ones).  Lines 1089-1092.
% ============================================================================
%% BLOCK 25 REPLACES:
%OLD% \emph{Numerical consistency} (deterministic Gil--Pelaez inversion, 511
%OLD% ordinates plus Gaussianized tail; computational, not part of the
%OLD% proof): at $m=0$, $1-\delta=4.073\times10^{-3}$ against the bounds
%OLD% $\min(0.156,\,e^{-1/(2C)}=0.0406)$; at $m=-0.1$,
%---NEW--->
\emph{Numerical consistency} ($\chi=\chi_4$; deterministic
Gil--Pelaez inversion, $511$
ordinates plus Gaussianized tail; computational, not part of the
proof): at $m=0$, $1-\delta=4.073\times10^{-3}$ against the bounds
$\min(0.156,\,e^{-1/(2C(\chi_4))}=0.0406)$; at $m=-0.1$,
%<---NEW---

% ============================================================================
% Proposition prop:unlogged restated for general chi.  Lines 1127-1138.
% ============================================================================
%% BLOCK 26 REPLACES:
%OLD% \begin{proposition}[transfer to the unlogged races]\label{prop:unlogged}
%OLD% Assume GRH$(\chi_4)$ and fix $m\in(-\tfrac12,\infty)$. Let
%OLD% $D_m(x)=\sum_{p\le x,\,p\equiv3(4)}p^m-\sum_{p\le x,\,p\equiv1(4)}p^m$
%OLD% and
%OLD% \[
%OLD% E_m'(x)\;=\;(2m+1)\,x^{-(m+1/2)}(\log x)\,D_m(x).
%OLD% \]
%OLD% Then $u\mapsto E_m'(e^u)$ has the \emph{same} limiting distribution
%OLD% $\nu_m$ as $E_m(e^u)$; consequently, under LI$(\chi_4)$, all of
%OLD% (ii)--(iii) hold verbatim for the unlogged race, with the same
%OLD% $\delta(m)$.
%OLD% \end{proposition}
%---NEW--->
\begin{proposition}[transfer to the unlogged races]\label{prop:unlogged}
Assume GRH$(\chi)$ and NCZ$(\chi)$, and fix $m\in(-\tfrac12,\infty)$.
Let
$D_m(x)=-\sum_{p\le x}\chi(p)p^m$ (for $\chi=\chi_4$ this is
$\sum_{p\le x,\,p\equiv3(4)}p^m-\sum_{p\le x,\,p\equiv1(4)}p^m$)
and
\[
E_m'(x)\;=\;(2m+1)\,x^{-(m+1/2)}(\log x)\,D_m(x).
\]
Then $u\mapsto E_m'(e^u)$ has the \emph{same} limiting distribution
$\nu_m$ as $E_m(e^u)$; consequently, under LI$(\chi)$, all of
(ii)--(iii) hold verbatim for the unlogged race, with the same
$\delta_\chi(m)$.
\end{proposition}
%<---NEW---

% ============================================================================
% Proof of prop:unlogged: |s| ~ |gamma| uses gamma_1(chi) > 0 (from
% NCZ), with chi-dependent constants.  Lines 1158-1159.
% ============================================================================
%% BLOCK 27 REPLACES:
%OLD% Write $s=\rho+m=m+\tfrac12+i\gamma$, so $|s|\asymp\gamma$ and
%OLD% $\mathrm{Re}\,s=m+\tfrac12>0$.
%---NEW--->
Write $s=\rho+m=m+\tfrac12+i\gamma$, so
$|s|\asymp_{m,\chi}|\gamma|$ (here the lower bound
$|\gamma|\ge\gamma_1(\chi)>0$ supplied by NCZ enters) and
$\mathrm{Re}\,s=m+\tfrac12>0$.
%<---NEW---

% ============================================================================
% Closing remark of the section: C -> C(chi); the m=0 numerical
% comparison is the chi_4 one.  Lines 1246-1256.
% ============================================================================
%% BLOCK 28 REPLACES:
%OLD% $\lim_Y\frac1Y\int_0^YE_m(e^u)^2\,du=1+\tfrac12\sum_\gamma|r_\gamma|^2
%OLD% =1+\sum_\gamma2a_\gamma(m)^2\le 1+C(2m+1)^2$, so the spread about the
%OLD% unit bias collapses quadratically in $(2m+1)$ under GRH alone. We do
%OLD% not pursue LI-free density statements (for which see \cite{Hay25},
%OLD% and \cite{Dev20} for the possible-atom phenomenology), because
%OLD% without LI $\nu_m$ could in principle carry an atom at $0$ and the
%OLD% density interpretation requires care. (2) The bounds of (iii) hold with
%OLD% $\sigma_m^2=\sum_\gamma 2a_\gamma(m)^2$ in place of $C(2m+1)^2$
%OLD% throughout, which is numerically slightly sharper (at $m=0$:
%OLD% $0.15557$ vs $0.15603$); we state the $C(2m+1)^2$ form because the
%OLD% $m$-dependence is transparent.
%---NEW--->
$\lim_Y\frac1Y\int_0^YE_m(e^u)^2\,du=1+\tfrac12\sum_\gamma|r_\gamma|^2
=1+\sum_\gamma2a_\gamma(m)^2\le 1+C(\chi)(2m+1)^2$, so the spread
about the
unit bias collapses quadratically in $(2m+1)$ under GRH alone. We do
not pursue LI-free density statements (for which see \cite{Hay25},
and \cite{Dev20} for the possible-atom phenomenology), because
without LI $\nu_m$ could in principle carry an atom at $0$ and the
density interpretation requires care. (2) The bounds of (iii) hold with
$\sigma_m^2=\sum_\gamma 2a_\gamma(m)^2$ in place of $C(\chi)(2m+1)^2$
throughout, which is numerically slightly sharper (for $\chi_4$ at
$m=0$:
$0.15557$ vs $0.15603$); we state the $C(\chi)(2m+1)^2$ form because
the
$m$-dependence is transparent.
%<---NEW---

% ============================================================================
% Monotonicity subsection: the Chernoff bound is no longer 'stated for
% chi_4' -- the theorem is general; point to the mod-3 corollary.
% Lines 1293-1296.
% ============================================================================
%% BLOCK 29 REPLACES:
%OLD% control we have is the Chernoff bound of Theorem~\ref{thm:interp}(iii)
%OLD% (stated for $\chi_4$; the same argument applies verbatim mod $3$), a
%OLD% monotone upper envelope for $1-\delta$ which is consistent with ---
%OLD% but does not imply --- monotonicity of $\delta$ itself.
%---NEW--->
control we have is the Chernoff bound of Theorem~\ref{thm:interp}(iii)
(applied with $\chi=\chi_4$, resp.\ $\chi_3$; cf.\
Corollary~\ref{cor:mod3rate}), a
monotone upper envelope for $1-\delta$ which is consistent with ---
but does not imply --- monotonicity of $\delta$ itself.
%<---NEW---

% ============================================================================
% Section 3 opening: the constant of thm:interp is now C(chi); this
% section evaluates the q=4 instance.  Lines 1337-1338.
% ============================================================================
%% BLOCK 30 REPLACES:
%OLD% consequences: the constant $C$ of Theorem~\ref{thm:interp} acquires a
%OLD% closed form (Corollary~\ref{cor:Cexact});
%---NEW--->
consequences: the constant $C=C(\chi_4)$ of Theorem~\ref{thm:interp}
(this section computes at the modulus $4$; the general identity is
\eqref{eq:varidq}) acquires a
closed form (Corollary~\ref{cor:Cexact});
%<---NEW---

% ============================================================================
% Section 4 preamble: thm:interp(ii) now identifies delta_q(m) with the
% race density for every real primitive chi, not only q=4.
% Lines 1731-1738.
% ============================================================================
%% BLOCK 31 REPLACES:
%OLD% For $q=4$, Theorem~\ref{thm:interp}(ii) identifies $\delta_4(m)$, under
%OLD% LI$(\chi_4)$, with the logarithmic density $\delta(m)$ of the set where
%OLD% the weighted race $R_m$ is positive; for other moduli $\delta_q(m)$ is,
%OLD% in this paper, a statement about the limiting random variable only (see
%OLD% the Scope remark at the end of the section). Note that the theorem
%OLD% below does \emph{not} assume LI: independence of the phases is built
%OLD% into the definition of $X_{m,\chi}$, and LI is what transfers the
%OLD% conclusion to the race.
%---NEW--->
For every real primitive $\chi$ under the standing hypotheses,
Theorem~\ref{thm:interp}(ii) identifies $\delta_q(m)$, under
LI$(\chi)$, with the logarithmic density $\delta_\chi(m)$ of the set
where
the weighted race $R_{m,\chi}$ is positive ($q=4$ and $q=3$ are the
computed cases). Note that the theorem
below does \emph{not} assume LI: independence of the phases is built
into the definition of $X_{m,\chi}$, and LI is what transfers the
conclusion to the race.
%<---NEW---

% ============================================================================
% Scope remark item (1) of the dissolution section: the race
% interpretation now follows for every real primitive chi from the
% generalized thm:interp(ii).  Lines 2359-2368.
% ============================================================================
%% BLOCK 32 REPLACES:
%OLD% (1) \emph{Race interpretation.} For $q=4$,
%OLD% Corollary~\ref{cor:dissolution-race} makes the dissolution law a
%OLD% statement about the weighted prime race itself. For $q=3$ the proof
%OLD% of Theorem~\ref{thm:interp} goes through verbatim (nothing in
%OLD% Lemmas~\ref{lem:trunc}--\ref{lem:squares} or the
%OLD% \cite{ANS14} verification is specific to the modulus $4$ beyond
%OLD% $\chi$ being real, primitive and nonprincipal), but we have not
%OLD% restated it, so for $q=3$ --- and for every other real primitive
%OLD% character --- Theorem~\ref{thm:dissolution} is, in this paper, a
%OLD% theorem about the limiting random variable $X_{m,\chi}$ only. (2)
%---NEW--->
(1) \emph{Race interpretation.} Theorem~\ref{thm:interp} is stated
and proved for every real primitive $\chi$ (under GRH$(\chi)$ and
$L(\tfrac12,\chi)\ne0$), so under the additional hypothesis LI$(\chi)$
its part (ii) transfers Theorem~\ref{thm:dissolution} to the weighted
prime race itself for \emph{any} such $\chi$;
Corollary~\ref{cor:dissolution-race} records the case $q=4$, and the
mod-$3$ densities of Table~\ref{tab:delta} instantiate $q=3$. Without
LI$(\chi)$, Theorem~\ref{thm:dissolution} is, as before, a
theorem about the limiting random variable $X_{m,\chi}$ only. (2)
%<---NEW---

% ============================================================================
% Methods, 'General real primitive modulus': no longer a record of an
% unproved general form -- the theorem is now stated and proved
% generally in Section 2; keep the what-changes summary.
% Lines 2615-2629.
% ============================================================================
%% BLOCK 33 REPLACES:
%OLD% \emph{General real primitive modulus.} For use beyond $q\in\{3,4\}$ we
%OLD% record the general form of Theorem~\ref{thm:interp} and of the variance
%OLD% identity. Let $\chi$ be a real primitive character mod $q$ with
%OLD% $\chi(-1)=(-1)^a$, $a\in\{0,1\}$, and assume GRH$(\chi)$, LI$(\chi)$,
%OLD% and the no-central-zero hypothesis
%OLD% \[
%OLD% \mathrm{NCZ}(\chi):\qquad L(\tfrac12,\chi)\neq0 .
%OLD% \]
%OLD% Define $R_m(x)=-\sum_{p\le x}\chi(p)p^m\log p$ (nonresidues minus
%OLD% residues), $E_m=(2m+1)x^{-(m+1/2)}R_m$, and
%OLD% $a_\gamma(m)=2M/\sqrt{M^2+\gamma^2}$ over the positive ordinates of
%OLD% $L(s,\chi)$. Then parts (i)--(iii) of Theorem~\ref{thm:interp} hold
%OLD% verbatim with $C(\chi)=\sum_{\gamma_\chi>0}2\gamma_\chi^{-2}$, and the
%OLD% variance identity becomes, with
%OLD% $\xi(s,\chi)=(q/\pi)^{(s+a)/2}\Gamma\bigl(\tfrac{s+a}2\bigr)L(s,\chi)$,
%---NEW--->
\emph{General real primitive modulus.} Theorem~\ref{thm:interp} is
stated and proved in Section~\ref{sec:theorem} for an arbitrary real
primitive character $\chi$ mod $q$ with
$\chi(-1)=(-1)^a$, $a\in\{0,1\}$, under GRH$(\chi)$, LI$(\chi)$,
and the no-central-zero hypothesis
\[
\mathrm{NCZ}(\chi):\qquad L(\tfrac12,\chi)\neq0
\]
(there $R_m(x)=-\sum_{p\le x}\chi(p)p^m\log p$, nonresidues minus
residues; $E_m=(2m+1)x^{-(m+1/2)}R_m$;
$a_\gamma(m)=2M/\sqrt{M^2+\gamma^2}$ over the positive ordinates of
$L(s,\chi)$; and parts (i)--(iii) hold with
$C(\chi)=\sum_{\gamma_\chi>0}2\gamma_\chi^{-2}$). Here we record what
the generalization changes relative to the mod-$4$ case, and the
general form of the variance identity: with
$\xi(s,\chi)=(q/\pi)^{(s+a)/2}\Gamma\bigl(\tfrac{s+a}2\bigr)L(s,\chi)$,
%<---NEW---

% ============================================================================
% Methods, same paragraph: the p | q prime-square terms are O_{q,m}(1)
% constants, not O_m(x^m log x) (the latter is false for m < 0, where
% x^m log x -> 0 while the terms are fixed positive constants).
% Lines 2635-2641.
% ============================================================================
%% BLOCK 34 REPLACES:
%OLD% What changes and what does not. \emph{Unchanged:} the prime-square bias
%OLD% term --- the $k=2$ prime powers contribute
%OLD% $\sum_{p\le\sqrt x,\,p\nmid q}p^{2m}\log p=x^{m+1/2}/(2m+1)+o(x^{m+1/2})$
%OLD% because $\chi(p)^2=1$ for every $p\nmid q$ and the finitely many $p\mid q$
%OLD% contribute $O_m(x^{m}\log x)$, so the normalized mean is $1$ for every
%OLD% real $\chi$; the Chernoff/Chebyshev bounds of (iii), which use only the
%OLD% arcsine moment generating function. \emph{Changed:} the gamma factor
%---NEW--->
What changes and what does not. \emph{Unchanged:} the prime-square bias
term --- the $k=2$ prime powers contribute
$\sum_{p\le\sqrt x,\,p\nmid q}p^{2m}\log p=x^{m+1/2}/(2m+1)+o(x^{m+1/2})$
because $\chi(p)^2=1$ for every $p\nmid q$, and the at most $\omega(q)$
primes $p\mid q$
contribute $O_{q,m}(1)$ in total (each term is
$p^{2m}\log p\le(1+q^{2m})\log q$, a constant; see the proof of
Lemma~\ref{lem:squares}), so the normalized mean is $1$ for every
real $\chi$; the Chernoff/Chebyshev bounds of (iii), which use only the
arcsine moment generating function. \emph{Changed:} the gamma factor
%<---NEW---

% ============================================================================
% Assumptions and limitations: the race transfer of the dissolution law
% holds for every real primitive chi via thm:interp(ii).
% Lines 2730-2735.
% ============================================================================
%% BLOCK 35 REPLACES:
%OLD% $2\max K+1$ in the continuity argument). The dissolution law,
%OLD% stated in the introduction as a Gaussian-comparison prediction, is
%OLD% proved with an error term in
%OLD% Theorem~\ref{thm:dissolution} (for the density of the limiting
%OLD% variable; the transfer to the race itself uses LI and is stated for
%OLD% $q=4$). The finite-$x$ log-measure for $(q,m)=(3,2)$ shows the
%---NEW--->
$2\max K+1$ in the continuity argument). The dissolution law,
stated in the introduction as a Gaussian-comparison prediction, is
proved with an error term in
Theorem~\ref{thm:dissolution} (for the density of the limiting
variable; the transfer to the race itself uses LI$(\chi)$ and holds
for every real primitive $\chi$ via Theorem~\ref{thm:interp}(ii),
Corollary~\ref{cor:dissolution-race} recording the case $q=4$). The
finite-$x$ log-measure for $(q,m)=(3,2)$ shows the
%<---NEW---

% ============================================================================
% REFEREE ADDENDUM.  Proof of lem:trunc, opening sentence: the claim
% "psi(t,chi)=0 for t<3" is chi_4-specific (chi_4(2)=0) and FALSE for
% any chi with 2 nmid q -- e.g. psi(2,chi_3) = chi_3(2) log 2 =
% -log 2 != 0, and chi_3 is now a headline case of the theorem.  The
% displayed identity itself remains exact for every chi (psi(t,chi)=0
% for t<2, and the Stieltjes lower limit 2^- captures the jump at
% t=2), so only the justification is corrected.  Line 635.
% ============================================================================
%% BLOCK 36 REPLACES:
%OLD% Since $\psi(t,\chi)=0$ for $t<3$, Riemann--Stieltjes partial summation
%OLD% gives exactly
%---NEW--->
Since $\psi(t,\chi)=0$ for $t<2$ (when $2\nmid q$ the term $n=2$ is
present, and the Stieltjes lower limit $2^-$ captures its jump),
Riemann--Stieltjes partial summation
gives exactly
%<---NEW---

% ============================================================================
% REFEREE ADDENDUM.  Step 2 of the proof of lem:classical: "absolute
% $A_1$" contradicts Block 7 ("all implied constants depend at most on
% $q$") -- $A_1$ comes from (F3), which is now $\ll_q$.
% Lines 556-557.
% ============================================================================
%% BLOCK 37 REPLACES:
%OLD% $G=\{\gamma:\,T-1<\gamma<T+2\}$; by (F3), $\#G\le A_1\log(T+2)$ for an
%OLD% absolute $A_1$. With $r=1/(4A_1\log(T+2))$ the union of the intervals
%---NEW--->
$G=\{\gamma:\,T-1<\gamma<T+2\}$; by (F3), $\#G\le A_1\log(T+2)$ for
some $A_1=A_1(q)$. With $r=1/(4A_1\log(T+2))$ the union of the
intervals
%<---NEW---

% ============================================================================
% REFEREE ADDENDUM.  Abstract, item (iii): the LI transfer of the
% dissolution law is no longer confined to the mod-4 race -- under
% LI(chi), Theorem thm:interp(ii) (now general) transfers it for every
% real primitive chi (as the body's Blocks 23/32/35 state).  Lines
% 56-58.
% ============================================================================
%% BLOCK 38 REPLACES:
%OLD% character under GRH and $L(\tfrac12,\chi)\neq0$ --- under LI it applies
%OLD% to the logarithmic density of the mod-$4$ race itself --- and verified
%---NEW--->
character under GRH and $L(\tfrac12,\chi)\neq0$ --- under LI$(\chi)$
it applies
to the logarithmic density of the race itself for every such $\chi$
--- and verified
%<---NEW---

% ============================================================================
% REFEREE ADDENDUM.  Results item 3: same stale mod-4-only transfer
% statement; the general transfer is Theorem thm:interp(ii), with
% cor:dissolution-race recording q=4.  Lines 244-247.
% ============================================================================
%% BLOCK 39 REPLACES:
%OLD% limiting variable of any real primitive character under GRH and
%OLD% $L(\tfrac12,\chi)\ne0$, and under LI$(\chi_4)$ for the logarithmic
%OLD% density of the mod-$4$ race itself
%OLD% (Corollary~\ref{cor:dissolution-race}).
%---NEW--->
limiting variable of any real primitive character under GRH and
$L(\tfrac12,\chi)\ne0$, and under LI$(\chi)$ for the logarithmic
density of the race itself, by Theorem~\ref{thm:interp}(ii)
(Corollary~\ref{cor:dissolution-race} records the case $q=4$).
%<---NEW---

